\section{Αποτελέσματα μηχανικών μοντέλων}

Τα παρακάτω αποτελέσματα που δίνονται, είναι μόνο για το \texttt{Support Vector Regression} με την συνάρτηση \texttt{normalize} σαν scaler, 
Παρόλο που άλλα μοντέλα, όπως τα \texttt{Random Forest Regressor} και \texttt{Gradient \linebreak Boosting Regressor},μπορεί να έβγαλαν καλύτερα αποτελέσματα σε κάποιες περιπτώσεις.
θα διατηρήσουμε το \texttt{Support Vector Regression} ,που είχε τα πιο σταθερά καλά αποτελέσματα. 

Θα παραθέσουμε τα αποτελέσματα για όλα τα δεδομένα, όπως και για δύο ξεχωριστές υποπεριπτώσεις επιλογής των δεδομένων με βάση τις θέσεις ρύπων πιο κοντά στους αισθητήρες.
Δηλαδή, θα πάρουμε όλα τα πειράματα για την πιο κοντινή θέση ρύπου σε κάθε αισθητήρα. Θα πάρουμε και για τους τρεις αισθητήρες, οπότε θα έχουμε για κάθε αισθητήρα και μηχανικό μοντέλο 
μια θέση που βρίσκεται κάτω από μισό μέτρο σε απόσταση και δύο ακόμα που θα είναι πάνω από ένα μέτρο.
Αντίστοιχα, θα μελετήσουμε την υποπερίπτωση όπου παίρνουμε τις δύο πιο κοντινές θέσης πηγή ρύπου ανά αισθητήρα, με παρόμοιο τρόπο όπως περιεγράφηκε για την πιο κοντινή θέση ρύπου.

Επίσης, θα πάρουμε σαν έξτρα υποπεριπτώσεις την χρήση ή όχι συνάρτησης διαχείρισης των outliers,μία περίπτωση για διαχείρσηση outliers στο σύνολο των δεδομένων και ένα ακόμα ανά στήλη.

% Θα τα τοποθετήσουμε στην αρχή 3 πίνακες,που παρουσιάζουν τα τελικά αποτελέσματα όλων των περιπτώσεων 
% Πιο κάτω θα παραθέσουμε τα διαγράμματα για την κάθε υποπερίπτωση,όπως και τους hyperparameters των τελικών μοντέλων:


\subsection{Αποτελέσματα πινάκων}

\setlength{\tabcolsep}{4pt}

%------------------------------------------------
\begin{table}[h!]
\centering
\caption{Αποτελέσματα – Όλα τα δεδομένα}
\label{tab:all_data}
\begin{tabularx}{\textwidth}{|
    >{\raggedright\arraybackslash}p{3cm} |
    >{\raggedright\arraybackslash}p{3cm} |
    X | X | X | X |}
\hline
\toprule
\textbf{Διαχείριση Outliers} & 
\textbf{Μέθοδος} &
\textbf{$R^2$} &
\textbf{MSE} &
\textbf{Μέση Ευκλείδεια Απόσταση} &
\textbf{Σχετικό Διάγραμμα} \\
\midrule
Όχι        & Όλα τα δεδομένα & 0.294992  & 0.684208  & 1.003398 & \autoref{fig:all_data_} \\
Ολική      & Όλα τα δεδομένα & -0.073574 & 1.036164  & 1.293871 & \autoref{fig:all_data_Global_Outliers} \\
Ανά στήλη  & Όλα τα δεδομένα & 0.175861  & 0.810631  & 1.160190 & \autoref{fig:all_data_Outliers_Per_Column} \\
\bottomrule
\end{tabularx}
\end{table}

%------------------------------------------------
\begin{table}[h!]
\centering
\small
\caption{Αποτελέσματα – Μόνο η πιο κοντινή πηγή}
\label{tab:closest_data}
\begin{tabularx}{\textwidth}{|
    >{\raggedright\arraybackslash}p{3cm} |
    >{\raggedright\arraybackslash}p{3cm} |
    X | X | X | X |}
\hline
\toprule
\textbf{Διαχείριση Outliers} & 
\textbf{Μέθοδος} &
\textbf{$R^2$} &
\textbf{MSE} &
\textbf{Μέση Ευκλείδεια Απόσταση} &
\textbf{Σχετικό Διάγραμμα} \\
\midrule
Όχι        & Πιο κοντινή πηγή & 0.636293 & 0.347179 & 0.699629 & \autoref{fig:closest} \\
Ολική      & Πιο κοντινή πηγή & 0.796589 & 0.350702 & 0.796589 & \autoref{fig:closest_Global_Outliers} \\
Ανά στήλη  & Πιο κοντινή πηγή & 0.835465 & 0.148242 & 0.519056 & \autoref{fig:closest_Outliers_Per_Column} \\
\bottomrule
\end{tabularx}
\end{table}

%------------------------------------------------
\begin{table}[h!]
\centering
\caption{Αποτελέσματα – Δύο πιο κοντινές πηγές}
\label{tab:second_closest_data}
\begin{tabularx}{\textwidth}{|
    >{\raggedright\arraybackslash}p{3cm} |
    >{\raggedright\arraybackslash}p{3cm} |
    X | X | X | X |}
\hline
\toprule
\textbf{Διαχείριση Outliers} & 
\textbf{Μέθοδος} &
\textbf{$R^2$} &
\textbf{MSE} &
\textbf{Μέση Ευκλείδεια Απόσταση} &
\textbf{Σχετικό Διάγραμμα} \\
\midrule
Όχι        & Δύο κοντινές πηγές & -0.421104 & 1.506079 & 1.469430 & \autoref{fig:two_closest} \\
Ολική      & Δύο κοντινές πηγές & 0.507745  & 0.572824 & 1.026803 & \autoref{fig:two_closest_Global_Outliers} \\
Ανά στήλη  & Δύο κοντινές πηγές & 0.481814  & 0.564319 & 1.002284 & \autoref{fig:two_closest_Outliers_Per_Column} \\
\bottomrule
\end{tabularx}
\end{table}